\documentclass[12pt,a4paper,oneside]{report}
\usepackage[margin=2.5cm]{geometry}
\usepackage{amsmath,amssymb,graphicx,booktabs,hyperref,xcolor,setspace,float,caption}
\onehalfspacing
\graphicspath{{../paper/figures/}}
\newcommand{\resultsabstract}{GIST reaches a spoofing recall of -- at 1\% false-positive rate on real 2026 traffic, where an integrity-threshold monitor reaches 0.01, and localizes simulated jammers among benign low-integrity aircraft with a median error of --\,km versus 42.5\,km for non-linear least squares.}
\newcommand{\datasetsentence}{We use 5 days (1st of each month, January 2024 to October 2026), 67\,M airborne trace points inside the study regions and 1.22\,M aircraft-windows.}
\newcommand{\splitsentence}{This gives 2 training, 1 validation and 2 test days.}
\newcommand{\simsentence}{Each training day contributes simulated scenes alongside real scenes.}

\newcommand{\todo}[1]{\textcolor{red}{[TODO: #1]}}

\begin{document}
\begin{titlepage}\centering
{\Large <<COLLEGE>>\par}\vspace{1cm}
{\large Final Year Project Report\par}\vspace{2cm}
{\huge\bfseries GIST: Detecting, Classifying and Localizing GNSS Jamming and Spoofing from Crowdsourced ADS-B with Deep Learning\par}\vspace{2cm}
{\large Submitted by\\ \textbf{<<AUTHOR>>} (<<ROLL\_NO>>)\par}\vspace{1cm}
{\large Under the guidance of\\ \textbf{<<GUIDE>>}\par}\vfill
{\large <<DEPARTMENT>>\\ <<MONTH YEAR>>\par}
\end{titlepage}

\chapter*{Certificate}
\todo{College certificate text, signed by the guide and head of department.}

\chapter*{Acknowledgements}
\todo{Acknowledgements.} The ADS-B data used in this work is provided by ADSB.lol and its contributors under the Open Database License.

\chapter*{Abstract}
GNSS (GPS) interference affects thousands of flights every day in several regions of the world. Two different threats are
involved: \emph{jamming}, which denies the satellite signal, and \emph{spoofing}, which replaces it with a false one.
Public monitoring maps threshold the Navigation Integrity Category (NIC) broadcast by aircraft; this detects jamming but
misses spoofing, since a spoofed receiver reports high integrity. This project builds GIST, a deep learning system that
treats every aircraft as a remote GNSS sensor. It derives physics-consistency features that a spoofer cannot fake cheaply,
builds weak labels from labeling functions and a generative label model so that no manual annotation is needed, trains a
trajectory Transformer with spatial graph attention to classify each aircraft-window as normal, jammed or spoofed, and
localizes jammers with a differentiable radio-propagation layer trained on simulated emitters injected into real traffic.
\resultsabstract

\tableofcontents
\listoffigures
\listoftables

\chapter{Introduction}
\section{Motivation}
Since 2022, interference with GNSS signals has become routine around the Baltic Sea, the Eastern Mediterranean, the Black
Sea, Iraq and Iran, Myanmar and the India--Pakistan border. Aircraft lose their satellite position, or worse, are fed a
confident but false one. Air navigation service providers, airlines and regulators need to know where interference happens,
whether it is jamming or spoofing, and where it comes from.
\section{Problem statement}
Given crowdsourced ADS-B broadcasts from all aircraft in a region, (1) detect GNSS interference at the level of individual
aircraft and 10-minute windows, (2) classify it as jamming or spoofing, and (3) estimate the location of jammers, without
manually labelled training data.
\section{Objectives}
\begin{enumerate}
\item Build an open, reproducible data pipeline from public ADS-B archives.
\item Design features that remain informative when the GNSS solution itself is manipulated.
\item Produce training labels by weak supervision.
\item Train a deep model that uses both each aircraft's history and its neighbours.
\item Localize jammers and compare with published classical methods.
\item Provide an interactive map for inspection.
\end{enumerate}
\section{Organisation of the report}
Chapter~2 gives background on GNSS and ADS-B. Chapter~3 reviews related work. Chapter~4 describes the data and weak
supervision, Chapter~5 the model, Chapter~6 the simulator, Chapter~7 the implementation, Chapter~8 the results and
Chapter~9 the challenges met. Chapter~10 concludes.

\chapter{Background}
\section{GNSS and its vulnerabilities}
GNSS receivers compute position from signals that arrive at the Earth's surface at around $-160$\,dBW, far below the
thermal noise floor. A transmitter of a few watts can therefore overwhelm the signal over hundreds of kilometres
(jamming). A spoofer broadcasts counterfeit signals that the receiver locks onto, producing a wrong position or time.
\section{ADS-B}
Automatic Dependent Surveillance--Broadcast (ADS-B) is the system by which aircraft broadcast their GNSS-derived position,
velocity and identity about twice per second on 1090\,MHz. Messages also carry integrity indicators: the Navigation Integrity
Category (NIC), Navigation Accuracy Category for position (NAC\textsubscript{p}), Source Integrity Level (SIL) and Geometric
Vertical Accuracy (GVA). Aircraft also report barometric altitude, which does not depend on GNSS, and geometric (GNSS)
altitude. Volunteer receiver networks such as adsb.lol and the OpenSky Network collect these broadcasts worldwide.
Multilateration (MLAT) between ground receivers gives an independent position for aircraft seen by several receivers.
\section{Why integrity thresholds miss spoofing}
Under jamming the receiver loses track and NIC falls towards zero. Under spoofing the receiver tracks a consistent but false
signal and continues to report a high NIC. Detecting spoofing therefore requires checking the reported position against
information the spoofer does not control.
\section{Deep learning components}
We use the Transformer encoder~\cite{vaswani2017attention} for per-aircraft sequences and attention over nearby aircraft
in the style of graph attention networks~\cite{brody2022gatv2}. Weak supervision follows the data-programming
approach~\cite{ratner2017snorkel} with a Dawid--Skene label model~\cite{dawid1979}.

\chapter{Literature Review}
Liu et al.~\cite{liu2022itm,liu2025nls} locate jammers with ADS-B by modelling received power with the Friis equation and
solving a non-linear least squares problem; on the Denver 2022 event they report about 4\,km error. Sch\"afer et
al.~\cite{jammerhorizon} use radio horizons to bound the jammer position. gpsjam.org~\cite{gpsjam} publishes daily maps of the
share of aircraft with low NIC per hexagon. The DLR study~\cite{eichstaedt2026sharc} measures global interference from
space-based ADS-B, and an industry white paper~\cite{aireon2026spoofing} describes spoofing checks based on integrity and
position consistency. None of these combine learned jam/spoof classification, cross-aircraft evidence and learned
localization with weak supervision, which is the gap this project addresses.

\chapter{Data and Weak Supervision}
\section{Data source and licence}
The adsb.lol \texttt{globe\_history} archive~\cite{adsblol} publishes one tarball of readsb trace files per day under ODbL-1.0.
ADS-B Exchange sample data was considered and rejected because its terms limit it to evaluation. \datasetsentence
\section{Processing}
Trace points are parsed in parallel and filtered to six interference regions and three comparison regions. Integrity fields,
which appear only in periodic detail objects, are forward-filled with their age recorded. Points are binned into 30-second steps
and 10-minute windows.
\section{Features}
\todo{Copy and extend Section III-B of the paper; include the wind-estimation fix from Chapter 9.}
\section{Labeling functions and the label model}
\begin{table}[t]\centering\small
\caption{Labeling functions (window = 10 min). ``modern'' = ADS-B version $\ge1$; ``cruise'' = median altitude $>10\,000$\,ft.}
\label{tab:lfs}
\begin{tabular}{@{}lll@{}}\toprule
LF & Vote & Condition \\\midrule
nic\_zero & JAM & modern, cruise, $\ge30\%$ steps NIC $\le1$ \\
nic\_drop & JAM & modern, cruise, NIC $\ge8$ seen and $\ge30\%$ steps NIC $\le6$ \\
pos\_loss & JAM & cruise, ADS-B gap $\ge90$\,s, $\ge10\%$ steps NIC $\le6$ \\
cell\_jam & JAM & $\ge30\%$ of cell aircraft and $\ge30\%$ own steps NIC $\le6$ \\
teleport & SPOOF & position jump $>10$\,km with NIC $\ge7$ \\
gdiff & SPOOF & $\ge2$ steps with $|\Delta h_{geo-baro}-\text{nbr}|>1000$\,ft, NIC $\ge7$ \\
cluster & SPOOF & member of a co-located, altitude-spread cluster \\
mlat$^\dagger$ & SPOOF & MLAT vs.\ ADS-B position disagreement $>10$\,km \\
clean & NORMAL & all NIC $\ge8$, no jump/gdiff/cluster, $\ge10$ steps \\
clean\_lowcell & NORMAL & modern, cell NIC-low share $<5\%$, own NIC never low \\\bottomrule
\multicolumn{3}{l}{$^\dagger$held out of training labels; evaluation only.}
\end{tabular}
\end{table}

\todo{Describe the label model, the all-abstain issue and its fix, and per-region label statistics.}

\chapter{The GIST Model}
\todo{Extend Section IV of the paper with an architecture diagram and the loss functions.}

\chapter{Interference Simulator}
\todo{Extend Section V of the paper with the propagation model, antenna pattern and the EIRP calibration.}

\chapter{Implementation}
\section{Software stack}
Python 3.12, PyTorch with the Apple MPS backend, Polars and DuckDB for data processing, H3 for spatial indexing, NumPy and
SciPy for the baselines. All code is in a single repository with tests for propagation and the label model.
\section{Hardware constraints}
All experiments ran on an Apple M3 laptop with 8\,GB of memory. This drove several design choices: streaming parse with raw
data deleted after processing, a 21$\times$21 localization grid with coarse-to-fine refinement instead of a dense grid,
padding scenes to four fixed sizes to avoid repeated GPU graph compilation, and capping the GPU memory cache.

\chapter{Results}
\begin{table*}[t]\centering\small
\caption{Test split (2026 days). Detection metrics per aircraft-window: macro-F1 over 3 classes, average precision (AP) and spoofing recall at 1\% false-positive rate (R@1\%). Real = label-model targets; Sim = exact simulated truth. Single-score baselines cannot separate jamming from spoofing; their anomaly score is used for both. Mean$\pm$std over seeds.}
\label{tab:results}
\begin{tabular}{@{}lcccc|cccc@{}}\toprule
 & \multicolumn{4}{c|}{Real traffic} & \multicolumn{4}{c}{Simulated emitters}\\
Method & F1 & Jam AP & Spoof AP & Spoof R@1\% & F1 & Jam AP & Spoof AP & Spoof R@1\% \\\midrule
gpsjam-style threshold & 0.62 & 0.86 & 0.01 & 0.01 & 0.26 & 0.93 & 0.07 & 0.01 \\
Isolation Forest & 0.56 & 0.80 & 0.04 & 0.06 & 0.27 & 0.48 & 0.15 & 0.01 \\
LSTM autoencoder & 0.38 & 0.37 & 0.03 & 0.06 & 0.27 & 0.58 & 0.13 & 0.03 \\
\textbf{GIST} (0 seeds) & -- & -- & -- & -- & -- & -- & -- & -- \\
\bottomrule\end{tabular}\vspace{6pt}

\begin{tabular}{@{}lcc|cc@{}}\toprule
 & \multicolumn{2}{c|}{Clean scenes} & \multicolumn{2}{c}{+5\% benign low-NIC aircraft}\\
Jammer localization [km] & median & p95 & median & p95 \\\midrule
Centroid of affected & 105.0 & 260.2 & 126.4 & 266.4 \\
Radio-horizon score & 121.3 & 382.3 & 138.0 & 448.6 \\
NLS (Liu et al.) & 26.0 & 301.7 & 42.5 & 315.4 \\
\textbf{GIST} & -- & -- & -- & -- \\
\bottomrule\end{tabular}
\end{table*}

\begin{table}[t]\centering\small
\caption{Ablations on the test split (seed 0 unless noted).}
\label{tab:ablations}
\begin{tabular}{@{}lccccc@{}}\toprule
 & \multicolumn{2}{c}{Real} & \multicolumn{2}{c}{Sim} & Loc. \\
Variant & F1 & Spoof R@1\% & F1 & Spoof R@1\% & med. km \\\midrule
GIST (full) & -- & -- & -- & -- & -- \\
w/o spatial graph attention & -- & -- & -- & -- & -- \\
integrity fields only & -- & -- & -- & -- & -- \\
w/o simulated training scenes & -- & -- & -- & -- & -- \\
\bottomrule\end{tabular}
\end{table}

\section{Case studies}
\todo{Insert case-study figures from paper/figures (case\_*.pdf) with discussion.}

\chapter{Challenges and Lessons}
\begin{itemize}
\item \textbf{Data licensing.} The first-choice dataset (ADS-B Exchange samples) was evaluation-only; switching to adsb.lol
(ODbL) kept the work publishable and gave daily rather than monthly coverage.
\item \textbf{Download throttling.} GitHub throttled release downloads to a few bytes per second after about 7\,GB; the
downloader was rewritten to resume per part and back off.
\item \textbf{Noisy wind feature.} A naive wind triangle from forward-filled airspeed and heading was useless (99th percentile
of several hundred knots in clean regions). Using readsb's own wind estimate only on points with fresh data, excluding
turns and comparing within altitude bands reduced the median deviation to about 5\,kt.
\item \textbf{Label-model artifact.} Windows on which every labeling function abstained were assigned to the class with the
highest abstention rate (spoofing), producing 13\% ``spoofed'' windows over Australia. These windows are now unlabelled.
\item \textbf{Unrealistic first simulator.} Emitters of 1--100\,W affected 0.4\% of aircraft per scene; raising the power to
the range estimated in the literature gave realistic affected regions.
\item \textbf{Memory on an 8\,GB laptop.} The training process grew to 7.6\,GB through the GPU caching allocator, slowing
epochs by an order of magnitude until the cache was capped.
\end{itemize}

\chapter{Conclusion and Future Work}
\todo{Summarise results.} Future work: the Denver 2022 case study with OpenSky data (pending data access), the MLAT-only
release as an independent position reference, receiver-aware radio-horizon checks, and an online version that updates
every few minutes.

\bibliographystyle{IEEEtran}
\bibliography{../paper/related}

\appendix
\chapter{Reproducing the results}
See \texttt{README.md}: fetch days, build labels, build scene datasets, train, evaluate, generate tables and figures.
\end{document}
